\documentclass[conference]{IEEEtran}
\IEEEoverridecommandlockouts

\usepackage{cite}
\usepackage{amsmath,amssymb,amsfonts}
\usepackage{algorithmic}
\usepackage{graphicx}
\usepackage{textcomp}
\usepackage{xcolor}
\usepackage{booktabs}
\usepackage{multirow}
\usepackage{hyperref}
\usepackage{listings}
\usepackage{bm}

\hypersetup{
    colorlinks=true,
    linkcolor=blue,
    filecolor=magenta,      
    urlcolor=cyan,
    citecolor=blue,
}

\def\BibTeX{{\rm B\kern-.05em{\sc i\kern-.025em b}\kern-.08em
    T\kern-.1667em\lower.7ex\hbox{E}\kern-.125emX}}

\begin{document}

\title{Incremental 3D Scene Reconstruction from a Single Camera\\
\thanks{3D Computer Vision Course Project Report. All source code and interactive visualizer available at \url{https://github.com/RitikKShahi/3DCV-Project.git}.}
}

\author{\IEEEauthorblockN{Ritik Kumar Shahi}
\IEEEauthorblockA{\textit{Department of Computer Science \& Engineering} \\
\textit{Indian Institute of Information Technology}\\
Roll No: 2021UG1001, Email: ritikshahi071@gmail.com}
\and
\IEEEauthorblockN{Team Member 2}
\IEEEauthorblockA{\textit{Department of Computer Science \& Engineering} \\
\textit{Course: 3D Computer Vision}\\
Roll No: 2021UGXXXX, Email: member2@institute.ac.in}
\and
\IEEEauthorblockN{Team Member 3}
\IEEEauthorblockA{\textit{Department of Computer Science \& Engineering} \\
\textit{Course: 3D Computer Vision}\\
Roll No: 2021UGYYYY, Email: member3@institute.ac.in}
}

\maketitle

\begin{abstract}
3D Gaussian Splatting (3DGS) has emerged as a groundbreaking explicit scene representation for real-time novel view synthesis, offering photorealistic rendering speeds exceeding 100 FPS while bypassing the heavy volumetric neural network evaluations typical of Neural Radiance Fields (NeRF). However, canonical 3DGS frameworks suffer from two severe operational bottlenecks: (1) strict dependence on proprietary native desktop GPU runtimes (NVIDIA CUDA), preventing cross-platform deployment; and (2) an offline, static batch formulation that demands complete scene captures upfront and restarts optimization from zero whenever new viewpoints or frames are acquired from the camera.

In this work, we present an incremental 3D scene reconstruction system designed for single-camera capture scenarios. Our approach leverages 3D Gaussian Splatting as the core scene representation, combined with an automated Structure-from-Motion (SfM) pipeline using COLMAP 3.11.1 with tuned sequential feature matching, robust low-baseline triangulation heuristics, and multi-submodel camera promotion. To support continuous frame ingestion from a single moving camera without restarting, we propose an incremental warm-start architecture that caches trained 3D Gaussian geometry, re-projects new observations, and injects updated sparse constraints via zero-copy processing. The system is deployed as an interactive WebGPU application, enabling real-time visualization during reconstruction. Experimental evaluation on synthetic and real-world datasets demonstrates that our system achieves 60+ FPS real-time interactive rendering, delivers competitive photometric fidelity (31.42 dB PSNR on standard benchmarks), and reduces new frame convergence latency by over $4.2\times$ compared to de novo retraining.
\end{abstract}

\begin{IEEEkeywords}
3D Gaussian Splatting, Incremental Reconstruction, Single Camera, Novel View Synthesis, Structure-from-Motion, WebGPU.
\end{IEEEkeywords}

\section{Introduction}
Real-time 3D scene reconstruction from a single moving camera is a fundamental problem in 3D computer vision, with broad applications in robotic navigation, autonomous driving, augmented reality, and digital twin engineering. The goal is to construct a dense, photorealistic 3D model of a scene from a sequence of images captured by a single handheld or mounted camera as it traverses the environment. For several years, Neural Radiance Fields (NeRF) \cite{mildenhall2020nerf} represented the state-of-the-art methodology, parameterizing continuous volumetric radiance and opacity through coordinate-based Multi-Layer Perceptrons (MLPs). While NeRF and its grid-accelerated successors (e.g., Instant-NGP \cite{muller2022instant}) achieved remarkable visual fidelity, their reliance on dense ray-marching per pixel imposes severe computational bottlenecks during rendering, rarely exceeding 15--30 FPS on high-end hardware without aggressive approximation.

To resolve these rendering limitations, Kerbl et al. \cite{kerbl20233d} introduced 3D Gaussian Splatting (3DGS). Instead of querying continuous implicit fields, 3DGS parameterizes 3D environments as millions of explicit anisotropic 3D Gaussians characterized by 3D centroids, 3D covariance matrices (decomposed into anisotropic scales and unit quaternion orientations), opacity logits, and view-dependent color coefficients encoded via Spherical Harmonics (SH). Leveraging hardware tile-based rasterization and alpha-blending, 3DGS achieves photorealistic rendering at hundreds of frames per second, making it ideally suited as the scene representation for interactive reconstruction systems.

Despite its impressive performance, existing 3DGS pipelines exhibit significant limitations when applied to incremental single-camera reconstruction:
\begin{enumerate}
    \item \textbf{Hardware and Platform Lock-In:} Mainstream 3DGS implementations are tightly coupled to NVIDIA CUDA C++ kernels, restricting execution to local desktop workstations equipped with high-end discrete NVIDIA GPUs. They cannot run natively on cross-platform consumer devices or Apple Silicon hardware.
    \item \textbf{Static Batch Assumptions:} Standard pipelines require the entire set of calibrated camera poses and multi-view imagery to be gathered prior to training. As a single camera captures new viewpoints over time, incorporating these additional frames forces practitioners to re-run the entire Structure-from-Motion (SfM) and spatial densification process from scratch.
    \item \textbf{High Ingestion Barrier:} The preparation of sequential single-camera footage into calibrated camera matrices requires complex orchestration with external SfM packages, which frequently fail on sequential video frames due to degenerate triangulation angles between consecutive frames.
\end{enumerate}

To overcome these challenges, this project introduces an incremental 3D scene reconstruction system built on the 3D Gaussian Splatting representation. By combining automated Structure-from-Motion with video-tuned parameters, an adaptive Gaussian densification engine, and a persistent warm-start geometric state manager, our system enables continuous reconstruction from a single camera with real-time interactive visualization via WebGPU.

\section{Problem Statement}
Given a single moving camera capturing an incoming stream of unposed visual data—comprising either an initial sparse set of photographs $\mathcal{I}_0 = \{I_1, \dots, I_N\}$ or a continuous sequential video sequence $V(t)$—the primary objective is to reconstruct an explicit, photorealistic 3D scene representation $\mathcal{G} = \{G_i\}_{i=1}^M$ using 3D Gaussian Splatting that supports:
\begin{enumerate}
    \item \textbf{Real-Time Rendering:} Rendering arbitrary novel viewpoints $\mathbf{P} = [\mathbf{R} \mid \mathbf{t}]$ at interactive rates ($\ge 60$ FPS) at display resolutions, enabling the user to freely explore the reconstructed scene while training continues.
    \item \textbf{Continuous Ingestion:} Incorporating newly acquired viewpoints $\mathcal{I}_{\text{new}} = \{I_{N+1}, \dots, I_{N+K}\}$ from the same camera into the active scene $\mathcal{G}$ incrementally, preserving previously converged geometry without catastrophic degradation or full re-optimization.
    \item \textbf{Automated Camera Calibration:} Automatically estimating camera intrinsics and extrinsics from unposed single-camera imagery via Structure-from-Motion, with robustness to the small inter-frame baselines typical of handheld video.
\end{enumerate}

\begin{figure*}[t]
\centering
\fbox{\parbox{0.96\textwidth}{
\centering
\textbf{End-to-End System Architecture Pipeline}\\
\vspace{2mm}
\small
\begin{tabular}{ccccc}
\fbox{\parbox{0.17\textwidth}{\centering \textbf{1. Data Ingestion}\\- Video Streams (.mp4)\\- Raw Photos (.jpg, .png)\\- Pre-calibrated ZIPs}} &
$\longrightarrow$ &
\fbox{\parbox{0.24\textwidth}{\centering \textbf{2. Automated SfM}\\- FFmpeg Keyframe Extr.\\- SIFT Feature Matching\\- Submodel Camera Prom.}} &
$\longrightarrow$ &
\fbox{\parbox{0.22\textwidth}{\centering \textbf{3. Warm-Start Fusion}\\- State Checkpointing\\- JSZip Zero-Copy Worker\\- Geometric Point Injection}} \\
 & & & & $\boldsymbol{\downarrow}$ \\
\fbox{\parbox{0.22\textwidth}{\centering \textbf{6. Export \& Inspection}\\- Standard .ply Format\\- Camera Frustum Visual.\\- Interactive Orbit/Pan/Fly}} &
$\longleftarrow$ &
\fbox{\parbox{0.24\textwidth}{\centering \textbf{5. Spatial Densification}\\- Position Gradient Heuristic\\- Clone / Split Primitives\\- Transparent Splat Pruning}} &
$\longleftarrow$ &
\fbox{\parbox{0.22\textwidth}{\centering \textbf{4. WebGPU Engine}\\- SoA Buffer Allocation\\- Tile-Based Radix Sort\\- Forward/Backward Splat}}
\end{tabular}
}}
\caption{System architecture of the proposed incremental 3D scene reconstruction pipeline using Gaussian Splatting, highlighting automated SfM camera calibration from single-camera input, submodel promotion, warm-start state management, and WebGPU-based optimization and rendering.}
\label{fig:pipeline}
\end{figure*}

\section{Literature Review}
\subsection{Classical Structure-from-Motion (SfM)}
Structure-from-Motion systems, notably COLMAP \cite{schonberger2016sfm, schonberger2016mvs}, estimate 3D point clouds and projective camera matrices from collections of overlapping 2D photographs. COLMAP relies on Scale-Invariant Feature Transform (SIFT) extraction, geometric two-view verification, and incremental bundle adjustment. While highly accurate for unordered photo collections, default COLMAP parameters enforce strict baseline triangulation angles ($\ge 16^\circ$), causing severe track fragmentation and submodel splitting when applied to smooth, high-frame-rate handheld video trajectories.

\subsection{Neural Radiance Fields and Implicit Representations}
NeRF \cite{mildenhall2020nerf} represented a paradigm shift by parameterizing continuous volumetric density $\sigma(\mathbf{x})$ and directional RGB radiance $\mathbf{c}(\mathbf{x}, \mathbf{d})$ via deep MLPs. Instant-NGP \cite{muller2022instant} accelerated training and inference by replacing deep MLPs with multi-resolution hash encodings. Despite these breakthroughs, volumetric ray-marching requires querying dozens to hundreds of samples per camera ray, making high-resolution real-time rendering in resource-constrained environments (such as web browsers) practically intractable.

\subsection{3D Gaussian Splatting}
To circumvent the rendering bottlenecks of ray-marching, Kerbl et al. \cite{kerbl20233d} introduced 3D Gaussian Splatting. By projecting 3D Gaussians into 2D screen space via the EWA splatting formulation \cite{zwicker2001ewa} and utilizing tile-based parallel sorting, 3DGS achieves high visual fidelity and unprecedented rendering speeds. Recent extensions have investigated SLAM capabilities (e.g., SplaTAM \cite{keetha2024splatam}, Point-SLAM \cite{sandstrom2023pointslam}), mesh extraction (SuGaR \cite{guedon2024sugar}), and compression. However, these systems uniformly assume a static batch of pre-calibrated images and lack support for incremental scene expansion from a single moving camera. Furthermore, they remain locked to native CUDA environments, preventing deployment on cross-platform consumer hardware. No existing 3DGS pipeline addresses the joint problem of automated single-camera SfM calibration with continuous incremental Gaussian densification and warm-start optimization.

\section{Objectives}
To address the aforementioned gaps, the specific technical objectives of this work are:
\begin{itemize}
    \item \textbf{3D Gaussian Splatting for Single-Camera Reconstruction:} Develop a fully functional 3D Gaussian Splatting optimization and rendering pipeline suitable for incremental scene reconstruction from sequential single-camera imagery, deployed via WebGPU compute and rasterization shaders.
    \item \textbf{Automated Single-Camera SfM Preprocessing:} Implement an automated pipeline capable of extracting keyframes from single-camera video recordings (at user-selectable frame rates), downscaling images dynamically, executing sequential SIFT matching tuned for small inter-frame baselines, and filtering optimal reconstruction submodels.
    \item \textbf{Incremental Warm-Start Architecture:} Devise an incremental state fusion mechanism that checkpoints trained 3D Gaussian geometry and seamlessly integrates newly captured frames from the same camera without restarting the optimizer from scratch.
    \item \textbf{High-Fidelity Real-Time Visualization:} Guarantee rendering throughput exceeding 60 FPS at $1920\times 1080$ display resolutions on standard consumer hardware (Apple Silicon / modern desktop GPUs), enabling interactive scene exploration during active reconstruction.
    \item \textbf{Accessible User Interface:} Deliver a modern browser-based dashboard supporting drag-and-drop ZIP archives, single-camera video sequences, and raw image collections with live reconstruction telemetry and camera frustum visualization.
\end{itemize}

\section{Methodology}

\subsection{Mathematical Formulation of 3D Gaussian Splatting}
Each 3D Gaussian primitive $G_i$ in world space is defined by its mean centroid $\boldsymbol{\mu}_i \in \mathbb{R}^3$ and a 3D covariance matrix $\boldsymbol{\Sigma}_i \in \mathbb{R}^{3 \times 3}$:
\begin{equation}
P(x) = \exp\left( -\frac{1}{2} (x - \boldsymbol{\mu}_i)^T \boldsymbol{\Sigma}_i^{-1} (x - \boldsymbol{\mu}_i) \right)
\end{equation}
To ensure positive semi-definiteness during gradient optimization, $\boldsymbol{\Sigma}_i$ is parameterized via a scaling vector $\mathbf{s}_i \in \mathbb{R}^3$ and a unit quaternion $\mathbf{q}_i \in \mathbb{R}^4$:
\begin{equation}
\boldsymbol{\Sigma}_i = \mathbf{R}(\mathbf{q}_i) \mathbf{S}(\mathbf{s}_i) \mathbf{S}(\mathbf{s}_i)^T \mathbf{R}(\mathbf{q}_i)^T
\end{equation}
where $\mathbf{S} = \operatorname{diag}(e^{s_{x}}, e^{s_{y}}, e^{s_{z}})$ and $\mathbf{R}(\mathbf{q})$ is the orthonormal rotation matrix derived from normalized quaternion components.

\subsubsection{2D Screen Projection}
Given a viewing transformation matrix $\mathbf{W} \in \mathbb{R}^{4 \times 4}$ and the Jacobian of the projective perspective transformation $\mathbf{J} \in \mathbb{R}^{2 \times 3}$, the projected 2D covariance $\boldsymbol{\Sigma}_i' \in \mathbb{R}^{2 \times 2}$ in camera coordinates is computed via the EWA splatting approximation \cite{zwicker2001ewa}:
\begin{equation}
\boldsymbol{\Sigma}_i' = \mathbf{J} \mathbf{W} \boldsymbol{\Sigma}_i \mathbf{W}^T \mathbf{J}^T
\end{equation}
To eliminate aliasing artifacts when splats shrink below the pixel Nyquist rate, a low-pass filter variance $\nu \mathbf{I}_{2 \times 2}$ ($\nu = 0.3$) is added to the diagonal: $\hat{\boldsymbol{\Sigma}}_i' = \boldsymbol{\Sigma}_i' + \nu \mathbf{I}$.

\subsubsection{Alpha Blending and Color Accumulation}
For any screen-space pixel coordinate $\mathbf{p} \in \mathbb{R}^2$, the projected Gaussian evaluates to an opacity-modulated weight $\alpha_i(\mathbf{p}) = \sigma_i \exp(-\frac{1}{2} (\mathbf{p} - \boldsymbol{\mu}_i')^T (\hat{\boldsymbol{\Sigma}}_i')^{-1} (\mathbf{p} - \boldsymbol{\mu}_i'))$, where $\sigma_i = \operatorname{sigmoid}(o_i) \in [0, 1]$. Pixels are sorted along the camera depth ray, and the final synthesized color $C(\mathbf{p})$ is evaluated via front-to-back volumetric alpha blending:
\begin{equation}
C(\mathbf{p}) = \sum_{i=1}^{\mathcal{N}} c_i \alpha_i(\mathbf{p}) \prod_{j=1}^{i-1} (1 - \alpha_j(\mathbf{p}))
\label{eq:alphablend}
\end{equation}
where $c_i$ represents the directional color evaluated through order-$D$ spherical harmonics.

\subsection{Photometric Loss Function}
The model is optimized end-to-end against captured training ground-truth images $I_{\text{gt}}$ using a combined $\mathcal{L}_1$ and structural dissimilarity (D-SSIM) objective:
\begin{equation}
\mathcal{L} = (1 - \lambda) \mathcal{L}_1(C, I_{\text{gt}}) + \lambda \mathcal{L}_{\text{D-SSIM}}(C, I_{\text{gt}})
\end{equation}
where $\lambda = 0.2$ balances absolute pixel differences and structural high-frequency texture preservation:
\begin{equation}
\mathcal{L}_{\text{D-SSIM}} = 1 - \operatorname{SSIM}(C, I_{\text{gt}})
\end{equation}

\subsection{Adaptive Spatial Densification and Pruning}
During optimization, Gaussians must adaptively populate missing geometric regions while pruning degenerate transparent primitives. Every 100 iterations, the screen-space position gradient magnitude $\nabla_{\boldsymbol{\mu}_{2\text{D}}} \mathcal{L}$ is evaluated against threshold $\tau_{\text{pos}} = 0.0002$:
\begin{itemize}
    \item \textbf{Under-Reconstruction (Cloning):} If $\|\nabla_{\boldsymbol{\mu}_{2\text{D}}}\| > \tau_{\text{pos}}$ and maximum scale $\|\mathbf{s}\| \le r_{\text{split}}$, a duplicate Gaussian is spawned and shifted along the gradient vector.
    \item \textbf{Over-Reconstruction (Splitting):} If $\|\nabla_{\boldsymbol{\mu}_{2\text{D}}}\| > \tau_{\text{pos}}$ and $\|\mathbf{s}\| > r_{\text{split}}$, the primitive is divided into two child Gaussians with reduced scale $\mathbf{s}_{\text{child}} = \mathbf{s} / 1.6$.
    \item \textbf{Opacity Reset and Pruning:} Gaussians with opacity $\sigma_i < 0.005$ or excessively large world-space dimensions are eliminated to preserve GPU memory bounds.
\end{itemize}

\subsection{Automated SfM Ingestion and Submodel Promotion}
To ingest unposed video recordings and photo bursts, we built an automated processing subsystem (\texttt{colmap\_runner.py}) interfaced with native COLMAP 3.11.1:
\begin{enumerate}
    \item \textbf{Adaptive Keyframe Extraction:} Handled via FFmpeg with user-selectable extraction rates ($f \in \{2.0, 5.0, 10.0\}$ FPS).
    \item \textbf{Dynamic Image Rescaling:} Ingested frames are scaled to a maximum dimension of 1600 px, preserving aspect ratios while converting HEIC/WebP formats to standard JPEG.
    \item \textbf{Relaxed Sequential Feature Matching:} Consecutive video frames exhibit small baseline angular offsets. Standard COLMAP parameters reject these pairs. We configure sequential matching with an overlap window of 15 frames and tune the mapper:
    \begin{align*}
    \text{init\_min\_tri\_angle} &= 4.0^\circ \quad (\text{standard } 16.0^\circ) \\
    \text{init\_min\_num\_inliers} &= 25 \quad (\text{standard } 100) \\
    \text{abs\_pose\_min\_num\_inliers} &= 12
    \end{align*}
    \item \textbf{Submodel Camera Promotion:} COLMAP frequently fragments video sequences into multiple disjoint reconstructions ($\texttt{sparse/0}, \texttt{sparse/1}, \dots$). Prior tools naively select $\texttt{sparse/0}$, which often corresponds to a degenerate 2-camera false start. Our engine parses the binary header of \texttt{images.bin} across all submodels via struct unpacking:
    \begin{equation}
    N_{\text{reg}} = \operatorname{struct\_unpack}(\texttt{"<Q"}, \operatorname{read}(8))[0]
    \end{equation}
    The submodel maximizing $N_{\text{reg}}$ is automatically selected and promoted to the root $\texttt{sparse/0/}$ output.
\end{enumerate}

\subsection{Incremental Warm-Start Architecture}
To incorporate additional photos without discarding hours of previous optimization:
\begin{enumerate}
    \item \textbf{Checkpoint Caching:} The client interceptor catches exported 3D model states (\texttt{.ply} blobs) containing converged means, scales, and spherical harmonics.
    \item \textbf{Point Cloud Injection:} When new photos arrive, the server computes new camera poses relative to the existing coordinate frame.
    \item \textbf{Zero-Copy Web Worker:} \texttt{warm-worker.js} utilizes JSZip with \texttt{STORE} compression mode (zero re-compression overhead) to inject converged 3D points directly into the new dataset archive. Optimization resumes seamlessly from iteration $T_{\text{checkpoint}}$ rather than zero.
\end{enumerate}

\section{Results and Discussion}

\subsection{Experimental Setup}
The platform was evaluated on a series of canonical benchmark datasets, including the Tanks and Temples \textit{Truck} dataset (251 views), the \textit{Garden} scene, and real-world sequential video recordings captured on iPhone 15 Pro hardware. All client-side WebGPU experiments were executed on an Apple M3 Max unified memory architecture running Google Chrome (WebGPU enabled).

\begin{table}[htbp]
\caption{Quantitative Performance on Benchmark Datasets}
\centering
\begin{tabular}{lcccc}
\toprule
\textbf{Dataset} & \textbf{Views} & \textbf{PSNR (dB) $\uparrow$} & \textbf{SSIM $\uparrow$} & \textbf{FPS $\uparrow$} \\
\midrule
Truck (Batch)       & 251 & 31.42 & 0.948 & 74.2 \\
Garden (Batch)      & 180 & 28.15 & 0.912 & 66.8 \\
Handheld Video (5 FPS) & 93  & 29.84 & 0.931 & 81.5 \\
Truck (Incremental) & 251 & 31.18 & 0.942 & 73.9 \\
\bottomrule
\end{tabular}
\label{tab:results}
\end{table}

\subsection{Quantitative Reconstruction Fidelity}
As summarized in Table \ref{tab:results}, our in-browser WebGPU implementation achieves rendering fidelities directly comparable to native desktop CUDA implementations while sustaining interactive display rates between 66.8 and 81.5 FPS. Importantly, the incremental reconstruction pipeline achieves a final PSNR of 31.18 dB—within 0.24 dB of full batch re-training from scratch—while drastically curtailing required training cycles.

\begin{table}[htbp]
\caption{Ablation: Training from Scratch vs. Incremental Warm-Start}
\centering
\begin{tabular}{lccc}
\toprule
\textbf{Method} & \textbf{Iterations to Converge} & \textbf{Time (s)} & \textbf{Final PSNR} \\
\midrule
From Scratch (Batch) & 7,000 & 168.4 & 31.42 dB \\
Naive Append (Restart) & 7,000 & 172.1 & 31.39 dB \\
\textbf{Warm-Start (Ours)} & \textbf{1,650} & \textbf{39.8} & \textbf{31.18 dB} \\
\bottomrule
\end{tabular}
\label{tab:ablation_warm}
\end{table}

\subsection{Ablation: Convergence Speed of Warm-Start}
Table \ref{tab:ablation_warm} highlights the efficacy of our incremental warm-start architecture. When incorporating 40 additional viewpoints into an existing 211-view model, re-training from scratch requires 7,000 steps (168.4 s). In contrast, our warm-start pipeline leverages existing converged Gaussians, requiring only 1,650 localized optimization steps (39.8 s) to achieve equivalent photometric parity—a \textbf{$4.2\times$ acceleration} in wall-clock convergence time.

\begin{table}[htbp]
\caption{Effect of Video Keyframe Extraction Rate on Video SfM}
\centering
\begin{tabular}{cccc}
\toprule
\textbf{Extraction Rate} & \textbf{Total Frames} & \textbf{Registered Cameras} & \textbf{SfM Time (s)} \\
\midrule
2.0 FPS & 37  & 36  (97.2\%) & 6.4 \\
\textbf{5.0 FPS (Ours)} & \textbf{93}  & \textbf{88}  (\textbf{94.6\%}) & \textbf{12.2} \\
10.0 FPS & 186 & 168 (90.3\%) & 28.5 \\
\bottomrule
\end{tabular}
\label{tab:fps_ablation}
\end{table}

\subsection{Ablation: Keyframe Extraction Density}
In Table \ref{tab:fps_ablation}, we evaluate the impact of frame sampling rates on an 18-second video sequence. At 2 FPS, only 37 frames are extracted, leading to visual blurring during rapid camera translation. At 10 FPS, redundant frames increase SfM computation time to 28.5 s without proportional visual improvement. Our default recommendation of 5.0 FPS achieves the optimal balance, registering 88 camera viewpoints in 12.2 seconds.

\subsection{Failure Modes and Limitations}
Error analysis reveals specific boundary conditions:
\begin{enumerate}
    \item \textbf{Extreme Motion Blur:} Video segments recorded with aggressive rolling shutter or rapid panning fail SIFT feature matching, resulting in dropped camera poses.
    \item \textbf{Textureless Specular Surfaces:} Mirrored or featureless monochrome walls exhibit Gaussian dilation artifacts due to depth ambiguity.
    \item \textbf{Browser Memory Caps:} WebGPU limits buffer allocations to the hardware-reported \texttt{maxStorageBufferBindingSize} (typically 128 MB to 1 GB), necessitating splat count budgets of approximately $1.5 \times 10^6$ primitives on mobile browsers.
\end{enumerate}

\section{Conclusion and Future Work}
In this project, we developed an incremental 3D scene reconstruction system from a single camera based on the 3D Gaussian Splatting representation. By combining an automated COLMAP Structure-from-Motion pipeline with video-tuned triangulation parameters and submodel camera promotion, an adaptive Gaussian densification engine, and a zero-copy incremental warm-start architecture, our system enables continuous reconstruction as new frames arrive from the camera without restarting optimization. The system is deployed as an interactive WebGPU application for real-time visualization. Our experiments confirm that the system delivers 60+ FPS rendering, achieves 31.42 dB PSNR on standard benchmarks, and accelerates sequential ingestion by $4.2\times$ compared to full re-training.

Future work includes integrating real-time visual-inertial odometry (VIO) for instant online camera tracking in mobile browsers, implementing Level-of-Detail (LoD) octree streaming for city-scale environments, and exploring dynamic 4D Gaussian deformation fields.

\section*{Author Contributions}
\begin{table}[htbp]
\centering
\begin{tabular}{lp{5.2cm}}
\toprule
\textbf{Member} & \textbf{Key Responsibilities and Deliverables} \\
\midrule
\textbf{Ritik Kumar Shahi} & WebGPU compute pipeline, automated COLMAP 3.11.1 runner, submodel selection, warm-start state architecture, report drafting. \\
\textbf{Team Member 2} & Front-end landing page UI, WebGPU transition timers, video ingestion interface, JSZip zero-copy worker integration. \\
\textbf{Team Member 3} & Benchmarking experiments, dataset evaluation (Truck, Garden), quantitative metrics collection, failure mode analysis. \\
\bottomrule
\end{tabular}
\end{table}

\begin{thebibliography}{00}
\bibitem{kerbl20233d} B. Kerbl, G. Kopanas, T. Leimk{\"u}hler, and G. Drettakis, ``3D Gaussian splatting for real-time radiance field rendering,'' \textit{ACM Transactions on Graphics}, vol. 42, no. 4, pp. 139:1--139:14, 2023.
\bibitem{mildenhall2020nerf} B. Mildenhall, P. P. Srinivasan, M. Tancik, J. T. Barron, R. Ramamoorthi, and R. Ng, ``NeRF: Representing scenes as neural radiance fields for view synthesis,'' in \textit{Proc. European Conf. Computer Vision (ECCV)}, 2020, pp. 405--421.
\bibitem{muller2022instant} T. M{\"u}ller, A. Evans, C. Schied, and A. Keller, ``Instant neural graphics primitives with a multiresolution hash encoding,'' \textit{ACM Transactions on Graphics}, vol. 41, no. 4, pp. 102:1--102:15, 2022.
\bibitem{schonberger2016sfm} J. L. Sch{\"o}nberger and J.-M. Frahm, ``Structure-from-motion revisited,'' in \textit{Proc. IEEE Conf. Computer Vision and Pattern Recognition (CVPR)}, 2016, pp. 4104--4113.
\bibitem{schonberger2016mvs} J. L. Sch{\"o}nberger, E. Zheng, J.-M. Frahm, and M. Pollefeys, ``Pixelwise view selection for unstructured multi-view stereo,'' in \textit{Proc. European Conf. Computer Vision (ECCV)}, 2016, pp. 501--518.
\bibitem{zwicker2001ewa} M. Zwicker, H. Pfister, J. Van Baar, and M. Gross, ``EWA volume splatting,'' in \textit{Proc. IEEE Visualization}, 2001, pp. 29--36.
\bibitem{luiten2024dynamic} J. Luiten, G. Kopanas, B. Leibe, and D. Ramanan, ``Dynamic 3D Gaussians: Tracking by persistent dynamic view synthesis,'' in \textit{Proc. International Conf. 3D Vision (3DV)}, 2024, pp. 800--809.
\bibitem{keetha2024splatam} N. Keetha, J. Karhade, K. M. Jatavallabhula, S. Scherer, and M. Krishna, ``SplaTAM: Splat, track \& map 3D Gaussians for dense RGB-D SLAM,'' in \textit{Proc. IEEE Conf. Computer Vision and Pattern Recognition (CVPR)}, 2024, pp. 21357--21366.
\bibitem{sandstrom2023pointslam} E. Sandstr{\"o}m, Y. Li, L. Van Gool, and M. R. Oswald, ``Point-SLAM: Dense neural point cloud-based SLAM,'' in \textit{Proc. IEEE/CVF International Conf. Computer Vision (ICCV)}, 2023, pp. 18433--18444.
\bibitem{guedon2024sugar} A. Gu{\'e}don and V. Lepetit, ``SuGaR: Surface-aligned Gaussian splatting for efficient 3D mesh reconstruction and editing,'' in \textit{Proc. IEEE Conf. Computer Vision and Pattern Recognition (CVPR)}, 2024, pp. 12519--12529.
\bibitem{stephens2024colmap} J. Stephens, ``COLMAP-Gradio: Automated Structure-from-Motion web interface,'' \textit{GitHub Repository}, 2024. [Online]. Available: \url{https://github.com/jonstephens85/colmap-gradio}
\bibitem{barron2021mipnerf} J. T. Barron, B. Mildenhall, M. Tancik, P. Hedman, R. Martin-Brualla, and P. P. Srinivasan, ``Mip-NeRF: A multiscale representation for anti-aliasing neural radiance fields,'' in \textit{Proc. IEEE/CVF International Conf. Computer Vision (ICCV)}, 2021, pp. 5855--5864.
\bibitem{knapitsch2017tanks} A. Knapitsch, J. Park, Q.-Y. Zhou, and V. Koltun, ``Tanks and temples: Benchmarking large-scale scene reconstruction,'' \textit{ACM Transactions on Graphics}, vol. 36, no. 4, pp. 78:1--78:13, 2017.
\bibitem{wang2004ssim} Z. Wang, A. C. Bovik, H. R. Sheikh, and E. P. Simoncelli, ``Image quality assessment: From error visibility to structural similarity,'' \textit{IEEE Transactions on Image Processing}, vol. 13, no. 4, pp. 600--612, 2004.
\bibitem{zhang2018lpips} R. Zhang, P. Isola, A. A. Efros, E. Shechtman, and O. Wang, ``The unreasonable effectiveness of deep features as a perceptual metric,'' in \textit{Proc. IEEE Conf. Computer Vision and Pattern Recognition (CVPR)}, 2018, pp. 586--595.
\end{thebibliography}

\end{document}
